% !PRINTED 346
\weberpage{346}
%<<CONT
thus if $\xi_1 + \xi_2 + \xi_3$ and $\eta_1 + \eta_2 + \eta_3$ are congruent to zero,
then the same holds of the left sides of (7).

Now conversely, under every linear substitution of the form (6), three quantities
($\xi, \eta$) of a triple always pass again into a triple, and from this it can be
proved that every equation of degree 9 whose Galois group $P$ is contained in the
complete linear congruence group is a triple equation, provided the further condition
is added that by the group $P$ any two given roots pass into two other, likewise
arbitrarily given roots — or, what amounts to the same thing, if $P$ is
\gesperrt{doubly transitive}.

For if $P$ consists entirely of linear substitutions, then by the adjunction of two
arbitrary roots $x_1 = (\xi_1, \eta_1)$, $x_2 = (\xi_2, \eta_2)$, $P$ reduces to a
group $P_1$ which leaves unchanged not only $x_1, x_2$, but also the third root
$x_3 = (\xi_3, \eta_3)$ of the triple determined by $x_1, x_2$; consequently $x_3$
belongs to the new domain of rationality, and hence $x_3$ is rationally expressible
in terms of $x_1, x_2$ in the form

\weq{8}{x_3 = \Theta (x_1, x_2).}

Now if the group $P$ is doubly transitive, there is a permutation in $P$ by which
$x_1, x_2$ pass into two arbitrary other roots $y_1, y_2$, and by which consequently
$x_3$ also passes into the third root of the triple $y_1, y_2, y_3$; and this
permutation is applicable to (8). Hence also

\[
y_3 = \Theta (y_1, y_2).
\]

We state this as a theorem:

3. \gesperrt{Every equation of degree 9 whose group is linear
and doubly transitive is a triple equation.}

By the permutations of a doubly transitive linear group every triple can be carried
into every other, and indeed in arbitrary arrangement.

The hypothesis of double transitivity can also be expressed thus: that all the
permutations of $P$ which leave one root unchanged still connect the remaining roots
transitively with one another. If we denote by $P_0$ the group contained in $P$ by
whose permutations one of the roots $x_0$ remains unchanged, then $P_0$ must still be
transitive with respect to the other roots. But if the group $P$ is only simply
transitive,
%CONT>>
